\documentclass[12pt,a4paper]{article}

% ---- Packages ----
\usepackage[margin=1in]{geometry}
\usepackage{graphicx}
\usepackage{float}
\usepackage{amsmath}
\usepackage{booktabs}
\usepackage{caption}
\usepackage{listings}
\usepackage{xcolor}
\usepackage{hyperref}
\usepackage{enumitem}
\usepackage[heading]{ctex}

\hypersetup{
    colorlinks=true,
    linkcolor=blue,
    urlcolor=blue,
    citecolor=blue
}

% ---- MATLAB code listing style ----
\lstdefinestyle{matlabstyle}{
    language=Matlab,
    basicstyle=\ttfamily\small,
    keywordstyle=\color{blue},
    commentstyle=\color{teal},
    stringstyle=\color{magenta},
    numbers=left,
    numberstyle=\tiny\color{gray},
    stepnumber=1,
    numbersep=8pt,
    breaklines=true,
    frame=single,
    columns=fullflexible,
    showstringspaces=false,
    extendedchars=true,
    inputencoding=utf8
}

\title{\textbf{Engineering Electromagnetic Fields Experiment 3}\\[6pt]
       \large Magnetic Field Distribution of Two Circular Current Loops}
\author{HE Yanning \\ 12311512 \\[4pt]
        Department of Electrical and Electronic Engineering \\
        Southern University of Science and Technology}
\date{\today}

\begin{document}
\maketitle

% ======================================================================
\section{Experiment Task Description}
% ======================================================================

This experiment uses MATLAB to compute and visualise the magnetic field generated by two parallel circular current loops. The calculation is based on the Biot--Savart law:
\begin{equation}\label{eq:bsl}
    d\mathbf{H} = \frac{I}{4\pi}\,\frac{d\mathbf{L} \times \mathbf{R}}{R^{3}},
\end{equation}
where $\mathbf{H}$ is the magnetic field intensity, $I\,d\mathbf{L}$ is the current element vector, and $\mathbf{R}$ is the vector from the current element to the field point.

\begin{figure}[H]
    \centering
    \includegraphics[width=0.60\textwidth]{images/1.png}
    \caption{Illustration of the Biot--Savart law: magnetic field $d\mathbf{H}$ produced by a current element $I\,d\mathbf{L}$ at a field point $P$. (Source: experiment handout.)}
    \label{fig:bsl}
\end{figure}

\begin{figure}[H]
    \centering
    \includegraphics[width=0.60\textwidth]{images/2.png}
    \caption{Geometry of a circular current loop and its magnetic field on the $yz$-plane. (Source: experiment handout.)}
    \label{fig:geometry}
\end{figure}

Fig.~\ref{fig:bsl} illustrates the Biot--Savart law for a single current element. Fig.~\ref{fig:geometry} shows the geometry of a circular loop in the $xy$-plane. For a circular loop of radius $a$ with its centre at the origin, the magnetic field on the central axis can be derived analytically. However, for field points off the axis, a closed-form solution is difficult. We therefore adopt a numerical approach: discretise the current loop into many small current elements, compute each element's contribution using Eq.~\eqref{eq:bsl}, and sum them using the superposition principle.

\subsection{Experiment Configuration}

Two circular loops, each of radius $a = 2\,\mathrm{m}$ and carrying current $I = 500\,\mathrm{A}$, are placed parallel to the $xy$-plane with centres at $(0,0,-1)$ and $(0,0,1)$. The magnetic field is evaluated on the $yz$-plane ($x = 0$). Two cases are studied:

\begin{enumerate}[label=(\roman*)]
    \item \textbf{Case 1 --- Same current direction:} Both loops carry current in the same sense. When the loop separation equals the radius ($d = a = 2\,\mathrm{m}$), this is known as a Helmholtz coil. Its key property is a nearly uniform magnetic field in the central region between the two loops.
    \item \textbf{Case 2 --- Opposite current directions:} The upper loop's current direction is reversed. The superposition of opposing fields produces a qualitatively different pattern.
\end{enumerate}

For each case, the following plots are generated: (a)~3D mesh of $|\mathbf{H}|$ over the full window $y,z \in [-5,5]$; (b)~3D mesh zoomed to the region of interest $y \in [-2,2], z \in [-1,1]$; (c)~contour map; (d)~field-line (streamline) plot; (e)~normalised vector plot. All figure titles include the student name and ID (HE Yanning, 12311512).

% ======================================================================
\section{Numerical Method}
% ======================================================================

The computation proceeds as follows. Each circular loop is divided into $N = 200$ straight segments of equal angular increment $\Delta\theta = 2\pi/N$. For each segment $i$, its midpoint coordinates $(x_{c,i},\, y_{c,i})$ and vector length components $(\Delta L_{x,i},\, \Delta L_{y,i})$ are computed. The field is evaluated on a regular grid of $100 \times 100$ points covering $y \in [-5,5]$, $z \in [-5,5]$.

At each grid point, the contributions from all $2 \times N$ current elements are summed:
\begin{equation}\label{eq:disc}
    \mathbf{H}(0, y_j, z_k) \approx \frac{I}{4\pi} \sum_{\text{loop}=1}^{2} \sum_{i=1}^{N}
    \frac{\Delta\mathbf{L}_{i}^{(\text{loop})} \times \bigl(\mathbf{r}_{jk} - \mathbf{r}_{c,i}^{(\text{loop})}\bigr)}
         {\bigl|\mathbf{r}_{jk} - \mathbf{r}_{c,i}^{(\text{loop})}\bigr|^{3}}.
\end{equation}

Because the observation plane is $x = 0$, the $H_x$ component is zero by symmetry. The total magnitude is $|\mathbf{H}| = \sqrt{H_y^{2} + H_z^{2}}$.

The complete MATLAB code is consolidated in \texttt{lab3\_code.mlx} (see Appendix~\ref{app:code}), with extensive comments explaining each step.

% ======================================================================
\section{Case 1: Same Current Direction (Helmholtz Coil)}
% ======================================================================

In this configuration, both loops carry current in the same direction. The magnetic fields of the two loops reinforce each other in the central region.

\begin{figure}[H]
    \centering
    \includegraphics[width=0.84\textwidth]{images/case1_01_full_mesh.png}
    \caption{Case~1: $|\mathbf{H}|$ distribution on the $yz$-plane, full view ($y,z \in [-5,5]$).}
    \label{fig:c1_m1}
\end{figure}

\begin{figure}[H]
    \centering
    \includegraphics[width=0.84\textwidth]{images/case1_02_local_mesh.png}
    \caption{Case~1: $|\mathbf{H}|$ distribution in the region of interest ($y \in [-2,2], z \in [-1,1]$).}
    \label{fig:c1_m2}
\end{figure}

Figs.~\ref{fig:c1_m1} and~\ref{fig:c1_m2} show the field magnitude. Two peaks appear near $z = \pm 1$, corresponding to the positions of the current-carrying loops. Inside the red rectangular region ($y \in [-2,2], z \in [-1,1]$), the surface is relatively flat, indicating that the field magnitude does not change dramatically.

\begin{figure}[H]
    \centering
    \includegraphics[width=0.84\textwidth]{images/case1_03_contour.png}
    \caption{Case~1: Contour map of $|\mathbf{H}|$. The red rectangle marks the region of interest.}
    \label{fig:c1_c}
\end{figure}

Fig.~\ref{fig:c1_c} displays isomagnetic (equal-field) contours. Green circles mark the four points where the two loops intersect the $yz$-plane; the overlaid markers ($*$ and $+$) indicate current direction (out of and into the page, respectively). The contours are dense near the loops (rapid field variation) and sparse near the centre (slow variation), which confirms the field uniformity.

\begin{figure}[H]
    \centering
    \includegraphics[width=0.84\textwidth]{images/case1_04_streamline.png}
    \caption{Case~1: Magnetic field lines on the $yz$-plane.}
    \label{fig:c1_s}
\end{figure}

The streamline plot (Fig.~\ref{fig:c1_s}) displays the field-line topology. Lines emerge from the interior of each loop, pass through the central region roughly parallel to the $z$-axis, and return through the exterior. The pattern is symmetric about $z = 0$.

\begin{figure}[H]
    \centering
    \includegraphics[width=0.84\textwidth]{images/case1_05_vector.png}
    \caption{Case~1: Normalised magnetic field vectors. Arrow lengths are uniform; arrow direction shows the local field direction.}
    \label{fig:c1_v}
\end{figure}

The vector plot (Fig.~\ref{fig:c1_v}) shows the local field direction at each grid point. Inside the red rectangle, the vectors are predominantly vertical and point in the same direction, consistent with the expected behaviour of a Helmholtz coil.

\textbf{Summary of Case~1:} The Helmholtz-coil arrangement produces a strong, approximately uniform magnetic field in the space between the two loops. Both the mesh plot and the contour map confirm the slow spatial variation of $|\mathbf{H}|$ in the central region.

% ======================================================================
\section{Case 2: Opposite Current Directions}
% ======================================================================

In this configuration, the upper loop's current direction is reversed. The two fields now oppose each other in the central region.

\begin{figure}[H]
    \centering
    \includegraphics[width=0.84\textwidth]{images/case2_01_full_mesh.png}
    \caption{Case~2: $|\mathbf{H}|$ distribution on the $yz$-plane, full view.}
    \label{fig:c2_m1}
\end{figure}

\begin{figure}[H]
    \centering
    \includegraphics[width=0.84\textwidth]{images/case2_02_local_mesh.png}
    \caption{Case~2: $|\mathbf{H}|$ distribution in the region of interest. Note the weaker field near $z=0$ compared with Case~1.}
    \label{fig:c2_m2}
\end{figure}

Comparing Figs.~\ref{fig:c2_m1}--\ref{fig:c2_m2} with their Case~1 counterparts, the field in the central area is noticeably weaker. This is because the $z$-components from the two loops point in opposite directions near $z = 0$ and partially cancel.

\begin{figure}[H]
    \centering
    \includegraphics[width=0.84\textwidth]{images/case2_03_contour.png}
    \caption{Case~2: Contour map of $|\mathbf{H}|$. Note the different marker pattern ($*$ and $+$) compared with Case~1, reflecting the reversed current in the upper loop.}
    \label{fig:c2_c}
\end{figure}

The contour map (Fig.~\ref{fig:c2_c}) shows a different marker arrangement from Case~1, reflecting the reversed current direction. The contours near the centre are more widely spaced and pulled outward, indicating a weaker and less uniform field.

\begin{figure}[H]
    \centering
    \includegraphics[width=0.84\textwidth]{images/case2_04_streamline.png}
    \caption{Case~2: Magnetic field lines on the $yz$-plane. The pattern is more complex than Case~1.}
    \label{fig:c2_s}
\end{figure}

The streamline plot (Fig.~\ref{fig:c2_s}) reveals a more complex field topology. Instead of passing straight through the central region, field lines form closed loops that directly connect the upper and lower coils.

\begin{figure}[H]
    \centering
    \includegraphics[width=0.84\textwidth]{images/case2_05_vector.png}
    \caption{Case~2: Normalised magnetic field vectors. Note the significant horizontal components near the centre.}
    \label{fig:c2_v}
\end{figure}

The vector plot (Fig.~\ref{fig:c2_v}) confirms that the field direction in the central region is no longer predominantly vertical. Substantial horizontal ($y$) components appear, and the vector orientation varies considerably from point to point.

\textbf{Summary of Case~2:} The opposite-current configuration destroys the field uniformity. The central field is weaker and exhibits significant directional variation. This highlights the importance of current direction in multi-coil electromagnetic systems.

% ======================================================================
\section{Discussion and Conclusion}
% ======================================================================

\subsection{Physical Interpretation}

The contrasting behaviours can be understood from the superposition principle. When currents flow in the same direction (Case~1), the $z$-components from both loops add constructively, producing a strong axial field. The $y$-components are anti-symmetric about $z = 0$ and cancel at the mid-plane, leaving a clean, vertical field.

When the currents oppose each other (Case~2), the $z$-components cancel near the mid-plane, and the residual field is dominated by $y$-components that form the characteristic inter-loop connection pattern visible in the streamline plot.

\subsection{Helmholtz Coil}

A Helmholtz coil is a pair of identical coaxial circular coils separated by a distance equal to their radius ($d = a$). In this experiment, $a = 2\,\mathrm{m}$ and the loop separation is also $2\,\mathrm{m}$, satisfying the Helmholtz condition. The numerical results confirm that Case~1 produces a nearly uniform field in the central region, which is why Helmholtz coils are widely used in laboratory applications requiring a controllable, uniform magnetic field.

\subsection{Numerical Considerations}

Each loop is divided into 200 segments, giving an angular resolution of $1.8^{\circ}$. The $100 \times 100$ grid provides a spatial resolution of $0.1\,\mathrm{m}$ in both $y$ and $z$. These parameters produce smooth field plots while keeping the computation manageable. The total number of Biot--Savart evaluations is approximately $4 \times 10^{6}$, which MATLAB handles efficiently through vectorised array operations.

\subsection{Conclusion}

This experiment demonstrated:
\begin{enumerate}[label=(\roman*)]
    \item The Biot--Savart law can be implemented numerically by discretising a current loop and applying the superposition principle;
    \item The same-direction (Helmholtz) configuration produces a strong, uniform magnetic field between the coils;
    \item Reversing one coil's current direction fundamentally alters the field distribution, destroying the central uniformity and creating a complex vector pattern;
    \item MATLAB provides a convenient platform for electromagnetic field computation and visualisation.
\end{enumerate}

% ======================================================================
\newpage
\appendix
\section{MATLAB Source Code: \texttt{lab3\_code.mlx}}\label{app:code}

The complete MATLAB code (including all Chinese annotations) is provided as \texttt{lab3\_code.mlx}. The code is also available as a plain-text \texttt{lab3\_code.m} file. Both files are located in the \texttt{source/} directory.

\subsection*{Code Structure}
\begin{itemize}
    \item \textbf{Parameter setup} (lines 1--30): Define loop radius, current, grid size, segment count, and the student tag.
    \item \textbf{Grid and loop discretisation} (lines 32--55): Build the $100 \times 100$ mesh and divide each loop into 200 current elements.
    \item \textbf{Case 1} (lines 57--170): Compute and plot all five figures for the Helmholtz-coil configuration (same current direction).
    \item \textbf{Case 2} (lines 172--290): Compute and plot all five figures for the opposite-current configuration.
\end{itemize}

\subsection*{Listing (plain-text version)}
\lstinputlisting[style=matlabstyle]{source/lab3_code.m}

\end{document}
